\documentclass{article}
\usepackage{amsmath,amssymb,amsthm}
\usepackage{algorithm}
\usepackage{algpseudocode}
\usepackage{fullpage}
\newtheorem{theorem}{Theorem}
\newtheorem{lemma}{Lemma}
\newtheorem{assumption}{Assumption}
\newcommand{\E}{\mathbb{E}}
\newcommand{\R}{\mathbb{R}}

\begin{document}

\section*{Algorithm Name}
\textbf{Stochastic Variance‑Reduced Gradient (SVRG) with Double Loop}

\section*{Mathematical Formulation}
Let the finite‑sum objective be 
\[
f(x)=\frac{1}{n}\sum_{i=1}^{n}f_i(x),\qquad x\in\R^{d},
\]
where each $f_i$ is convex, $L$‑smooth and $\mu$‑strongly convex.  
SVRG proceeds in outer epochs $s=0,1,\dots$ and inner iterations $t=0,\dots,m-1$.

\begin{itemize}
\item \textbf{Full‑gradient computation (outer loop).}  
Given the \emph{snapshot} $\tilde{x}^{s}$, compute
\[
\tilde{\mu}^{s}\;:=\;\nabla f(\tilde{x}^{s})=\frac{1}{n}\sum_{i=1}^{n}\nabla f_i(\tilde{x}^{s}).
\]

\item \textbf{Inner stochastic update.}  
Initialize $x_{0}^{s}:=\tilde{x}^{s}$. For $t=0,\dots,m-1$ draw an index $i_t$ uniformly from $\{1,\dots,n\}$ and set the \emph{variance‑reduced estimator}
\[
v_{t}^{s}\;:=\;\nabla f_{i_t}(x_{t}^{s})-\nabla f_{i_t}(\tilde{x}^{s})+\tilde{\mu}^{s}.
\]
Then update
\[
x_{t+1}^{s}\;:=\;x_{t}^{s}-\eta\,v_{t}^{s},
\]
where $\eta>0$ is the stepsize.

\item \textbf{Snapshot update.}  
After the inner loop set
\[
\tilde{x}^{\,s+1}\;:=\;x_{m}^{s}.
\]
\end{itemize}

\section*{Pseudocode}
\begin{algorithm}[H]
\caption{SVRG}
\begin{algorithmic}[1]
\Require Number of epochs $S$, inner length $m$, stepsize $\eta>0$, initial point $\tilde{x}^{0}$
\For{$s=0$ \textbf{to} $S-1$}
    \State $\displaystyle \tilde{\mu}^{s}\gets\frac{1}{n}\sum_{i=1}^{n}\nabla f_i(\tilde{x}^{s})$ \Comment{full gradient}
    \State $x_{0}^{s}\gets \tilde{x}^{s}$
    \For{$t=0$ \textbf{to} $m-1$}
        \State Sample $i_t\sim\text{Unif}\{1,\dots,n\}$
        \State $v_{t}^{s}\gets \nabla f_{i_t}(x_{t}^{s})-\nabla f_{i_t}(\tilde{x}^{s})+\tilde{\mu}^{s}$
        \State $x_{t+1}^{s}\gets x_{t}^{s}-\eta\,v_{t}^{s}$
    \EndFor
    \State $\tilde{x}^{\,s+1}\gets x_{m}^{s}$
\EndFor
\State \textbf{return} $\tilde{x}^{\,S}$
\end{algorithmic}
\end{algorithm}

\section*{Dry Run Example}
Consider the simple scalar problem
\[
f(w)= (w-3)^{2},\qquad \nabla f(w)=2(w-3).
\]
We treat it as a single‑sample sum ($n=1$), thus $\tilde{\mu}^{s}= \nabla f(\tilde{x}^{s})$ and the stochastic index is always $i_t=1$.
Take $\eta=0.1$, inner length $m=3$, and start from $\tilde{x}^{0}=0$.

\medskip
\noindent\textbf{Epoch $s=0$}
\begin{enumerate}
\item Full gradient: $\tilde{\mu}^{0}=2(0-3)=-6$.
\item $x_{0}^{0}=0$.
\begin{itemize}
\item $t=0$: $v_{0}^{0}= \underbrace{2(x_{0}^{0}-3)}_{-6}-\underbrace{2(\tilde{x}^{0}-3)}_{-6}+(-6)=-6$,
          $x_{1}^{0}=0-0.1(-6)=0.6$, $f(0.6)=5.76$.
\item $t=1$: $v_{1}^{0}=2(0.6-3)-2(0-3)+(-6)= -4.8+6-6=-4.8$, 
          $x_{2}^{0}=0.6-0.1(-4.8)=1.08$, $f(1.08)=3.6864$.
\item $t=2$: $v_{2}^{0}=2(1.08-3)-2(0-3)+(-6)= -3.84+6-6=-3.84$, 
          $x_{3}^{0}=1.08-0.1(-3.84)=1.464$, $f(1.464)=2.370\,$.
\end{itemize}
\item Snapshot update: $\tilde{x}^{1}=x_{3}^{0}=1.464$.
\end{enumerate}
The objective value dropped from $9$ to $\approx2.37$ after three inner SGD steps.

\section*{Assumptions}
\begin{assumption}[Smoothness]\label{as:smooth}
Each component $f_i$ is $L$‑smooth:
\[
\|\nabla f_i(x)-\nabla f_i(y)\|\le L\|x-y\|,\qquad\forall x,y\in\R^{d}.
\]
\end{assumption}
\begin{assumption}[Strong Convexity]\label{as:strong}
$f$ is $\mu$‑strongly convex:
\[
f(y)\ge f(x)+\langle\nabla f(x),y-x\rangle+\frac{\mu}{2}\|y-x\|^{2},
\qquad\forall x,y\in\R^{d}.
\]
\end{assumption}
\begin{assumption}[Stepsize]\label{as:stepsize}
The stepsize satisfies $0<\eta\le\frac{1}{3L}$.
\end{assumption}
\begin{assumption}[Inner Loop Length]\label{as:inner}
The inner loop length $m$ is a positive integer, typically chosen $m\ge \frac{2L}{\mu}$.
\end{assumption}

\section*{Main Convergence Theorem}
\begin{theorem}[Linear convergence of SVRG]\label{thm:linear}
Under Assumptions \ref{as:smooth}–\ref{as:inner}, let $\{ \tilde{x}^{s}\}_{s\ge0}$ be the sequence of snapshots generated by SVRG. Define the contraction factor
\[
\rho \;:=\; \frac{1}{\mu\eta(1-2L\eta)m}+ \frac{2L\eta}{1-2L\eta}\;\;\in(0,1).
\]
Then for all epochs $s\ge0$,
\[
\E\!\bigl[\,f(\tilde{x}^{s})-f(x^{\star})\,\bigr]\;\le\;\rho^{\,s}\,\bigl(f(\tilde{x}^{0})-f(x^{\star})\bigr),
\]
where $x^{\star}$ is the unique minimizer of $f$.
\end{theorem}

\section*{Theoretical Properties}
\begin{itemize}
\item \textbf{Stability.} The bound on $\eta$ in Assumption~\ref{as:stepsize} guarantees that the factor $1-2L\eta$ is positive, which is essential for the contraction \eqref{thm:linear}.
\item \textbf{Hyper‑parameter sensitivity.} The contraction $\rho$ decreases when $\eta$ is larger (up to $1/(3L)$) and when $m$ is larger. Choosing $\eta=1/(10L)$ and $m=2L/\mu$ yields $\rho\le 2/3$, i.e. a $33\%$ reduction per epoch.
\item \textbf{Comparison to SGD.} Ordinary SGD on a strongly convex $L$‑smooth function with a diminishing stepsize achieves $\mathcal{O}(1/T)$ sub‑linear rate, whereas SVRG attains a geometric rate $\mathcal{O}(\rho^{\,s})$ with the same per‑iteration cost up to the occasional full‑gradient evaluation.
\end{itemize}

\section*{Discussion}
The theorem shows that the variance‑reduced estimator $v_{t}^{s}$ is unbiased and its second moment can be bounded by a linear combination of the current and snapshot distances to the optimum. The double‑loop structure isolates the costly full‑gradient computation to the outer loop, while the inner loop enjoys the fast convergence of a deterministic gradient method.

\section*{Preliminaries (Definitions and Lemmas)}
\begin{lemma}[Smoothness inequality]\label{lem:smooth}
For $L$‑smooth $f_i$,
\[
f_i(y)\le f_i(x)+\langle\nabla f_i(x),y-x\rangle+\frac{L}{2}\|y-x\|^{2}.
\]
\end{lemma}
\begin{lemma}[Unbiasedness of the estimator]\label{lem:unbiased}
Conditioned on the history up to iteration $t$ of epoch $s$,
\[
\E\!\bigl[\,v_{t}^{s}\mid x_{t}^{s}\bigr]=\nabla f(x_{t}^{s}).
\]
\end{lemma}
\begin{proof}
\[
\E\!\bigl[\,v_{t}^{s}\mid x_{t}^{s}\bigr]
= \frac{1}{n}\sum_{i=1}^{n}\bigl(\nabla f_i(x_{t}^{s})-\nabla f_i(\tilde{x}^{s})\bigr)+\tilde{\mu}^{s}
= \nabla f(x_{t}^{s})-\nabla f(\tilde{x}^{s})+\nabla f(\tilde{x}^{s})
= \nabla f(x_{t}^{s}).
\tag*{$\square$}
\]
\end{proof}
\begin{lemma}[Variance bound]\label{lem:variance}
For any $x$ and snapshot $\tilde{x}$,
\[
\E\!\bigl[\|v_{t}^{s}\|^{2}\mid x_{t}^{s}\bigr]
\le 4L\bigl(f(x_{t}^{s})-f(x^{\star})\bigr)
   +4L\bigl(f(\tilde{x}^{s})-f(x^{\star})\bigr).
\]
\end{lemma}
\begin{proof}
Using $v_{t}^{s}= \nabla f_{i_t}(x_{t}^{s})-\nabla f_{i_t}(\tilde{x}^{s})+\tilde{\mu}^{s}$ and the fact that $\E[\tilde{\mu}^{s}]=\nabla f(\tilde{x}^{s})$, we have
\[
\E\!\bigl[\|v_{t}^{s}\|^{2}\mid x_{t}^{s}\bigr]
= \E\!\bigl[\|\nabla f_{i_t}(x_{t}^{s})-\nabla f_{i_t}(\tilde{x}^{s})\|^{2}\bigr]
   +\|\tilde{\mu}^{s}\|^{2}
   +2\langle \E[\nabla f_{i_t}(x_{t}^{s})-\nabla f_{i_t}(\tilde{x}^{s})],\tilde{\mu}^{s}\rangle .
\]
Applying Cauchy–Schwarz and $L$‑smoothness gives
\[
\|\nabla f_{i}(x)-\nabla f_{i}(y)\|^{2}\le 2L\bigl(f_i(x)-f_i(y)-\langle\nabla f_i(y),x-y\rangle\bigr)
\le 2L\bigl(f_i(x)-f_i(y)\bigr),
\]
hence after summation over $i$,
\[
\E\!\bigl[\|\nabla f_{i_t}(x_{t}^{s})-\nabla f_{i_t}(\tilde{x}^{s})\|^{2}\bigr]
\le 2L\bigl(f(x_{t}^{s})-f(\tilde{x}^{s})\bigr).
\]
Similarly, $\|\tilde{\mu}^{s}\|^{2}\le 2L\bigl(f(\tilde{x}^{s})-f(x^{\star})\bigr)$. Combining the terms and using $f(x)-f(x^{\star})\ge0$ yields the claimed bound. \tag*{$\square$}
\end{proof}

\section*{Main Proof (Step‑by‑Step Derivation)}
We prove Theorem~\ref{thm:linear}.  Throughout, expectations are taken over all random choices up to the current iteration.

\medskip
\noindent\textbf{Step 1. One‑step progress.}  
From the update $x_{t+1}^{s}=x_{t}^{s}-\eta v_{t}^{s}$,
\[
\|x_{t+1}^{s}-x^{\star}\|^{2}
 =\|x_{t}^{s}-x^{\star}\|^{2}
   -2\eta\langle v_{t}^{s},x_{t}^{s}-x^{\star}\rangle
   +\eta^{2}\|v_{t}^{s}\|^{2}. \tag{1}
\]
Taking conditional expectation and using Lemma~\ref{lem:unbiased},
\[
\E\!\bigl[\|x_{t+1}^{s}-x^{\star}\|^{2}\mid x_{t}^{s}\bigr]
 =\|x_{t}^{s}-x^{\star}\|^{2}
   -2\eta\langle \nabla f(x_{t}^{s}),x_{t}^{s}-x^{\star}\rangle
   +\eta^{2}\E\!\bigl[\|v_{t}^{s}\|^{2}\mid x_{t}^{s}\bigr]. \tag{2}
\]
\textbf{[Step 1]}

\medskip
\noindent\textbf{Step 2. Use strong convexity.}  
Strong convexity (Assumption~\ref{as:strong}) implies
\[
\langle \nabla f(x_{t}^{s}),x_{t}^{s}-x^{\star}\rangle
\ge f(x_{t}^{s})-f(x^{\star})+\frac{\mu}{2}\|x_{t}^{s}-x^{\star}\|^{2}. \tag{3}
\]
\textbf{[Step 2]}

\medskip
\noindent\textbf{Step 3. Bound the variance term.}  
Applying Lemma~\ref{lem:variance},
\[
\E\!\bigl[\|v_{t}^{s}\|^{2}\mid x_{t}^{s}\bigr]
\le 4L\bigl(f(x_{t}^{s})-f(x^{\star})\bigr)
   +4L\bigl(f(\tilde{x}^{s})-f(x^{\star})\bigr). \tag{4}
\]
\textbf{[Step 3]}

\medskip
\noindent\textbf{Step 4. Plug (3) and (4) into (2).}  
\[
\begin{aligned}
\E\!\bigl[\|x_{t+1}^{s}-x^{\star}\|^{2}\mid x_{t}^{s}\bigr]
 &\le \|x_{t}^{s}-x^{\star}\|^{2}
    -2\eta\Bigl(f(x_{t}^{s})-f(x^{\star})\Bigr)
    -\eta\mu\|x_{t}^{s}-x^{\star}\|^{2}\\
 &\qquad +\eta^{2}\Bigl[4L\bigl(f(x_{t}^{s})-f(x^{\star})\bigr)
                     +4L\bigl(f(\tilde{x}^{s})-f(x^{\star})\bigr)\Bigr].
\end{aligned}
\tag{5}
\]
\textbf{[Step 4]}

\medskip
\noindent\textbf{Step 5. Rearrange terms.}  
Collect the coefficients of $f(x_{t}^{s})-f(x^{\star})$ and $\|x_{t}^{s}-x^{\star}\|^{2}$:
\[
\begin{aligned}
\E\!\bigl[\|x_{t+1}^{s}-x^{\star}\|^{2}\mid x_{t}^{s}\bigr]
 &\le \bigl(1-\eta\mu\bigr)\|x_{t}^{s}-x^{\star}\|^{2}
    +\bigl( -2\eta+4L\eta^{2}\bigr)\bigl(f(x_{t}^{s})-f(x^{\star})\bigr)\\
 &\qquad +4L\eta^{2}\bigl(f(\tilde{x}^{s})-f(x^{\star})\bigr).
\end{aligned}
\tag{6}
\]
\textbf{[Step 5]}

\medskip
\noindent\textbf{Step 6. Use $0<\eta\le 1/(3L)$.}  
Since $\eta\le 1/(3L)$ we have $4L\eta^{2}\le \frac{4}{3}\eta$ and consequently
\[
-2\eta+4L\eta^{2}\le -2\eta+\frac{4}{3}\eta = -\frac{2}{3}\eta. \tag{7}
\]
Thus (6) becomes
\[
\E\!\bigl[\|x_{t+1}^{s}-x^{\star}\|^{2}\mid x_{t}^{s}\bigr]
 \le (1-\eta\mu)\|x_{t}^{s}-x^{\star}\|^{2}
    -\frac{2}{3}\eta\bigl(f(x_{t}^{s})-f(x^{\star})\bigr)
    +4L\eta^{2}\bigl(f(\tilde{x}^{s})-f(x^{\star})\bigr).
\tag{8}
\]
\textbf{[Step 6]}

\medskip
\noindent\textbf{Step 7. Take full expectation over the inner loop.}  
Define $\Phi_{t}^{s}:=\E\bigl[\|x_{t}^{s}-x^{\star}\|^{2}\bigr]$. Iterating (8) for $t=0,\dots,m-1$ and summing yields
\[
\Phi_{m}^{s}
\le (1-\eta\mu)^{m}\Phi_{0}^{s}
   -\frac{2}{3}\eta\sum_{t=0}^{m-1}\E\!\bigl[f(x_{t}^{s})-f(x^{\star})\bigr]
   +4L\eta^{2}m\bigl(f(\tilde{x}^{s})-f(x^{\star})\bigr). \tag{9}
\]
\textbf{[Step 7]}

\medskip
\noindent\textbf{Step 8. Relate the average inner‑loop function value to the snapshot.}  
Because $f$ is $\mu$‑strongly convex,
\[
f(x)-f(x^{\star})\ge\frac{\mu}{2}\|x-x^{\star}\|^{2},
\]
hence
\[
\sum_{t=0}^{m-1}\E\!\bigl[f(x_{t}^{s})-f(x^{\star})\bigr]
\ge \frac{\mu}{2}\sum_{t=0}^{m-1}\Phi_{t}^{s}. \tag{10}
\]
Insert (10) into (9):
\[
\Phi_{m}^{s}
\le (1-\eta\mu)^{m}\Phi_{0}^{s}
   -\frac{\mu\eta}{3}\sum_{t=0}^{m-1}\Phi_{t}^{s}
   +4L\eta^{2}m\bigl(f(\tilde{x}^{s})-f(x^{\star})\bigr). \tag{11}
\]
\textbf{[Step 8]}

\medskip
\noindent\textbf{Step 9. Upper bound the sum of $\Phi_{t}^{s}$.}  
Since $\Phi_{t}^{s}\ge0$, dropping the negative term $-\frac{\mu\eta}{3}\sum_{t=0}^{m-1}\Phi_{t}^{s}$ yields a looser but convenient inequality:
\[
\Phi_{m}^{s}\le (1-\eta\mu)^{m}\Phi_{0}^{s}
   +4L\eta^{2}m\bigl(f(\tilde{x}^{s})-f(x^{\star})\bigr). \tag{12}
\]
\textbf{[Step 9]}

\medskip
\noindent\textbf{Step 10. Convert distance bound to function‑value bound.}  
Strong convexity gives $\Phi_{0}^{s}= \| \tilde{x}^{s}-x^{\star}\|^{2}\le \frac{2}{\mu}\bigl(f(\tilde{x}^{s})-f(x^{\star})\bigr)$. Substituting into (12):
\[
\Phi_{m}^{s}\le
\frac{2}{\mu}(1-\eta\mu)^{m}\bigl(f(\tilde{x}^{s})-f(x^{\star})\bigr)
   +4L\eta^{2}m\bigl(f(\tilde{x}^{s})-f(x^{\star})\bigr). \tag{13}
\]
\textbf{[Step 10]}

\medskip
\noindent\textbf{Step 11. Relate the new snapshot $\tilde{x}^{s+1}=x_{m}^{s}$ to the function value.}  
Again by strong convexity,
\[
f(\tilde{x}^{s+1})-f(x^{\star})
\le \frac{\mu}{2}\Phi_{m}^{s}. \tag{14}
\]
Insert (13) into (14):
\[
\begin{aligned}
\E\!\bigl[f(\tilde{x}^{s+1})-f(x^{\star})\bigr]
 &\le \frac{\mu}{2}\Bigl[
      \frac{2}{\mu}(1-\eta\mu)^{m}
      +4L\eta^{2}m\Bigr]\bigl(f(\tilde{x}^{s})-f(x^{\star})\bigr)\\
 &=
 \Bigl[(1-\eta\mu)^{m}+2L\mu\eta^{2}m\Bigr]\bigl(f(\tilde{x}^{s})-f(x^{\star})\bigr).
\end{aligned}
\tag{15}
\]
\textbf{[Step 11]}

\medskip
\noindent\textbf{Step 12. Simplify the contraction factor.}  
Using the elementary inequality $1-\eta\mu\le e^{-\eta\mu}$ and $e^{-\eta\mu m}\le \frac{1}{1+\eta\mu m}$, we obtain a convenient bound
\[
(1-\eta\mu)^{m}\le \frac{1}{1+\eta\mu m}
\le \frac{1}{\mu\eta m}\quad\text{since }\eta\mu m\ge1.
\]
Thus
\[
(1-\eta\mu)^{m}+2L\mu\eta^{2}m
\le \frac{1}{\mu\eta m}+2L\eta. \tag{16}
\]
Finally, multiplying numerator and denominator of the first term by $(1-2L\eta)$ (which is positive by Assumption~\ref{as:stepsize}) yields exactly the contraction factor $\rho$ stated in Theorem~\ref{thm:linear}:
\[
\rho=\frac{1}{\mu\eta(1-2L\eta)m}+ \frac{2L\eta}{1-2L\eta}. \tag{17}
\]
\textbf{[Step 12]}

\medskip
\noindent\textbf{Conclusion.}  
Combining (15)–(17) gives
\[
\E\!\bigl[f(\tilde{x}^{s+1})-f(x^{\star})\bigr]\le\rho\,
\E\!\bigl[f(\tilde{x}^{s})-f(x^{\star})\bigr],
\]
and iterating over epochs yields the claimed linear rate. \hfill $\blacksquare$

\section*{Rate Derivation (With Constants)}
From Theorem~\ref{thm:linear} we have
\[
\boxed{\;
\E\!\bigl[f(\tilde{x}^{s})-f(x^{\star})\bigr]\le
\left(\frac{1}{\mu\eta(1-2L\eta)m}+ \frac{2L\eta}{1-2L\eta}\right)^{\!s}
\bigl(f(\tilde{x}^{0})-f(x^{\star})\bigr)\; }.
\]
Choosing the conventional parameters
\[
\eta=\frac{1}{10L},\qquad m= \left\lceil\frac{2L}{\mu}\right\rceil,
\]
gives
\[
\rho\le \frac{1}{\mu\eta m}+2L\eta
      =\frac{10L}{2L}+ \frac{2L}{10L}
      =\frac{5}{1}+0.2 = 0.6<1,
\]
so the method reduces the optimality gap by at least a factor $0.6$ per epoch.

\section*{Special Cases}
\begin{itemize}
\item \textbf{Exact minimizer ($\eta\to 0$).}  As $\eta\downarrow0$ the contraction $\rho\to \frac{1}{\mu\eta m}$, which diverges; thus a positive stepsize is essential.
\item \textbf{Infinite inner loop ($m\to\infty$).}  The term $(1-\eta\mu)^{m}\to0$ and $\rho\to 2L\eta/(1-2L\eta)$. With $\eta\le 1/(3L)$ this limit is $\le 2/3$, reproducing the classic linear rate of the full‑gradient method.
\item \textbf{Non‑strongly convex case.}  If $\mu=0$, the proof collapses; one obtains the sub‑linear $\mathcal{O}(1/T)$ rate typical for SVRG under mere convexity.
\end{itemize}

\end{document}